% Generated by report_gen.py (hwde v1) - do not hand-edit.
\documentclass[11pt,a4paper]{article}
\usepackage[margin=2.2cm]{geometry}
\usepackage{graphicx}
\usepackage{booktabs}
\usepackage{longtable}
\usepackage{pdfpages}
\makeatletter
\@ifundefined{KV@Gin@artifact}{\define@key{Gin}{artifact}[]{}}{}
\makeatother
\usepackage{xcolor}
\usepackage[hidelinks]{hyperref}
\setcounter{tocdepth}{1}
\makeatletter
\def\ps@hwde{\let\@mkboth\@gobbletwo\def\@oddhead{}\def\@evenhead{}\def\@oddfoot{\small PCB-0018-A -- bldc-motor-driver\hfil\thepage}\let\@evenfoot\@oddfoot}
\makeatother
\pagestyle{hwde}
\begin{document}
\sloppy
\begin{center}
{\LARGE\bfseries bldc-motor-driver --- Full Design Document}\\[6pt]
{\Large\bfseries PCB-0018-A}\\[6pt]
{\large hwde v1 pipeline}\\[2pt]
generated 2026-09-28 08:39:24
\end{center}
\tableofcontents

\section{Board and Run Metadata}
{\small
\begin{longtable}{lp{11cm}}
\toprule
\textbf{Field} & \textbf{Value} \\
\midrule
\endhead
board & bldc-motor-driver \\
part number & PCB-0018-A \\
workspace & \texttt{PCB-0018-A\_bldc-motor-driver} \\
phase & P9 \\
gate status & all 5 recorded gates pass \\
state created & 2026-09-27T14:02:00 \\
state updated & 2026-09-28T08:38:51 \\
\bottomrule
\end{longtable}
}
\section{Overview}
\subsection*{PCB-0018-A bldc-motor-driver ? brief}
Owner brief (2026-09-27): "also do a three phase motor driver". The reference is a published demo board seen in a 15-second video; frame: \texttt{reference-video-frame.png} (owner's screenshot). No spec was published, so the spec below is designed from a general brief and informed by that frame. Every line marked (guess) is our own inference, not a published fact.

\subsubsection*{What the frame shows}
\begin{itemize}
\item Three half-bridges, Q1-Q6, in 5x6-class power packages; low-side shunts labelled SHUNT A/B/C.
\item A VM power rail with one large electrolytic (C11 area), VM and GND pour islands, wire pads/test points TP1-TP5
on the edges (large plated pads: we read them as VM, GND and the three phase outputs).
\item A QFP MCU (\textasciitilde{}48-64 pins) in the middle, a buck to V3V3 (L1, BUCK\_SW net), several SOIC/TSSOP ICs (U1-U10:
gate driver, current-sense amps, regulators), two tactile switches SW1/SW2, a small JST-style connector and a
wider white connector (FFC/header), mounting holes, many small diodes (D10-D23), roughly square, 4 layers.
\end{itemize}
\subsubsection*{Spec (ours)}
\begin{itemize}
\item Motor: 3-phase BLDC/PMSM, sensored or sensorless FOC. (guess)
\item Supply VM: 12-24 V nominal, 10-28 V operating, 40 V-rated power stage, reverse-polarity and transient
protection on VM (TVS). (guess)
\item Current: 10 A RMS per phase continuous at 25 C ambient without a heatsink is the target, 20 A peak; derate
shown by a thermal estimate. (guess)
\item Power stage: 6 N-channel MOSFETs, 40 V, low Rds(on), 5x6 mm power packages; bootstrapped high side.
PWM 20-30 kHz. Dead time set in hardware or the MCU.
\item Gate drive: a 3-phase or 3x half-bridge gate driver with bootstrap, at \textgreater{}= 10 V VGS across the whole VM range
(a gate-drive rail independent of VM if needed); gate resistors, turn-off path, and pull-downs so FETs stay
off with the MCU in reset.
\item Current sensing: three low-side shunts with Kelvin (4-wire) connections to differential current-sense
amplifiers, filtered into the MCU's ADCs; overcurrent comparator/fault to the timer break input.
\item Voltage sensing: VM divider and three phase-voltage dividers (BEMF) to MCU ADCs; an NTC at the power stage.
\item MCU: STM32G4 (motor-control timers, fast ADCs), LQFP-48 or -64. (guess: G431/G474 class)
\item Power: buck VM -\textgreater{} 3.3 V (and whatever the gate driver needs), rated for the MCU + sensors + headroom.
\item I/O: a small JST connector for control (UART/CAN or PWM/step input, guess), a wider connector for Hall/encoder
(5 V, GND, H1-H3 / A-B-Z), SWD programming, two buttons (reset + user), status LEDs.
\item Motor and power: VM, GND, PHASE A/B/C on large plated pads (wire solder or screw lugs).
\item 4-layer, SMD assembly at JLCPCB (THT only where a connector or the electrolytic needs it), silkscreen
"PCB-0018 REV A".
\item Correctness first: gate-drive loop, shunt Kelvin sensing, power commutation loop, thermal ? then size. Size is an
output: aim for about 60 x 60 mm or smaller, roughly square like the reference board.
\item Stop at a DRC/ERC-clean fab package. Do NOT order. End with COMPARISON.md against what the video frame shows.
\end{itemize}
\section{Requirements}
\subsection*{PCB-0018-A bldc-motor-driver ? requirements}
\subsubsection*{1. Function}
A three-phase BLDC/PMSM motor driver board, designed from a general spec and informed by one frame of a published 3-phase motor-driver demo board (no spec was published). It drives a 3-phase motor via six N-channel MOSFET half-bridges with a gate driver, senses per-phase current via low-side shunts and phase/ bus voltage via dividers, and closes the control loop on an onboard MCU (sensored or sensorless FOC). It exposes control I/O (UART/CAN or step/PWM), a Hall/encoder connector, SWD programming, two buttons, and status LEDs. No mode token is present in the brief; this is a normal (non-mode) run, scope and binding are the default P5 board\_init flow, size is stated as a target with an explicit "size is an output" note (see Section 5).

\subsubsection*{2. Interfaces}
\begin{itemize}
\item \textbf{Motor phase outputs}: PHASE A/B/C, large plated pads (wire solder or
screw lugs), rated to the 20 A peak power spec. ASSUMED: pad-only, no
connector shell (frame shows bare pads, not a housed connector).
\item \textbf{Power input}: VM + GND, large plated pads (same pad group as phase
outputs going by the frame's TP1-TP5 layout). ASSUMED: pad-only, no
connector shell.
\item \textbf{Control connector}: small JST-style connector. ASSUMED signal set:
UART (TX/RX) or CAN, plus a power/GND pair ? exact protocol is an open
question (Q1).
\item \textbf{Hall/encoder connector}: wider white connector (FFC or header).
ASSUMED pinout: 5V, GND, H1-H3 (or A/B/Z quadrature) ? exact pin count/
pitch is an open question (Q2).
\item \textbf{Programming}: SWD (2-4 pin header or pads), not visible as a distinct
connector in the frame ? ASSUMED present per brief text.
\item \textbf{User I/O}: two tactile switches (SW1 reset, SW2 user), status LED(s)
(not clearly resolved in the frame; ASSUMED at least 1 status LED near
the MCU).
\item \textbf{Test points}: TP1-TP5, large plated pads read as VM, GND, and the
three phase outputs (test/probe access, not a load-bearing connector).
\end{itemize}
\subsubsection*{3. Power}
\begin{itemize}
\item \textbf{Input rail VM}: PROVISIONAL 12-24 V nominal, 10-28 V operating range,
power stage rated to 40 V. These current/voltage figures are the
orchestrator's provisional answers, not this analyst's guesses ? see
Section 8 (safety flags) and Q3/Q4 for confirmation.
\item \textbf{Protection}: reverse-polarity and transient (TVS) protection on VM.
ASSUMED needed given "motor supply, external wiring" context; not
confirmed by the frame.
\item \textbf{Onboard rails}: buck VM -\textgreater{} 3.3V (L1/BUCK\_SW net visible in frame) for
MCU + digital sensors. A separate gate-drive rail (\textgreater{}=10V VGS, independent
of VM if VM can sag below 10V) is required ? GUESS: rail voltage and
whether it's a dedicated regulator or bootstrap-only is an open question
(Q5).
\item \textbf{Rail budget}: not computed at this stage (P0); flagged for P2/P3
power-tree design. GUESS budget only: MCU + gate driver + sense amps +
LEDs, low tens of mA on 3V3, gate-drive rail sized to switching losses
at 20-30 kHz.
\item \textbf{Battery/charging}: none. This is a bench/system supply-fed motor
driver, not a battery-powered device (frame shows no charge IC or
battery connector).
\end{itemize}
\subsubsection*{4. Environment}
Not stated in the brief or resolvable from the frame. ASSUMED (low risk, default): indoor/lab use, 0-40C ambient, no enclosure, no IP rating, no vibration spec ? this is a demo/prototype driver board, not a certified product. Thermal: no heatsink assumed per brief ("without a heatsink" is the stated target condition), so thermal derating from the 6 power FETs and shunts is a P2/P6 design concern, not a requirement input here.

\subsubsection*{5. Size \& mounting}
\begin{itemize}
\item Target outline: approximately 60mm x 60mm or smaller, roughly square,
matching the frame's proportions. Per the brief, {*}{*}size is explicitly an
output, not an input{*}{*}: "Correctness first ... then size." This is a
RELAXABLE preference, not a HARD cap ? no HARD cap is set. The frame's
own board is a plausible upper bound but is not to be treated as binding;
placement should optimize for the electrical/thermal correctness items
first (gate-drive loop, Kelvin shunt sensing, commutation loop, thermal)
and let final size follow.
\item Mounting holes: 8 holes visible in the frame (H1-H10 labels seen, several
are likely test points not holes ? frame shows large filled circles at
roughly the 4 corners plus 4 edge-midpoints = up to 8 mounting holes).
ASSUMED: standard M3 mounting holes, corner + edge pattern, exact count
and diameter TBD at layout (P5) ? not a hard requirement to pin down at
P0.
\item Height limit: none stated; no enclosure implied. The single large
electrolytic capacitor is the tallest component visible; no clearance
constraint given.
\item Layer count: 4 layers, stated directly in the brief (HARD ? this is a
stated spec, not geometry-under-optimization).
\end{itemize}
\subsubsection*{6. Quantity \& budget}
Not stated in the brief. ASSUMED (default, low risk): prototype quantity, 1-5 boards, no target unit cost set ? this is a reverse-engineering/demo exercise, not a cost-optimized production run. Open question Q6 offers a default.

\subsubsection*{7. Assembly}
JLC PCBA, SMD assembly, THT only where a connector or the electrolytic capacitor needs it (stated directly in the brief). Single-sided placement of fine-pitch parts assumed typical (double-sided PCBA not called out; ASSUMED single-sided SMD placement, with THT parts on the primary side) ? open question Q7 offers a default.

\subsubsection*{8. Compliance/safety flags}
This board triggers multiple safety-relevant flags. Per the pipeline rule, none of these are guessed ? they are recorded as open questions or as PROVISIONAL values pending orchestrator/owner confirmation:

\begin{itemize}
\item \textbf{Motor drive}: yes ? 3-phase BLDC/PMSM motor driver, spinning mass
hazard under fault or miswiring.
\item \textbf{High current (\textgreater{}3 A)}: yes ? PROVISIONAL 10 A RMS per phase continuous,
20 A peak (orchestrator-provided, not this analyst's guess). This is a
hard-confirm item: see Q3.
Bulk-cap derating (P4, 2026-09-27): bus-cap RMS ripple is \textasciitilde{}0.5-0.6x phase
RMS, so 10 A/phase puts 5-6 A into the bus caps and most of it into the
electrolytics at 20 kHz. Two 470 uF/63 V cans (C101, C104) share it, but
each is rated about 1-2 A (a guess: the part's datasheet was not
available), so 10 A continuous is above what they are rated for. Raising
PWM toward 40 kHz moves ripple into the MLCC bank. The P8 thermal check
must state the continuous rating the caps allow.
\item \textbf{Voltage up to 28 V}: yes ? PROVISIONAL VM 10-28 V operating, 40 V
power-stage rating (orchestrator-provided). Not mains voltage (\textless{}50V), so
no mains isolation/creepage-for-mains requirement applies, but 28 V at
20 A peak still needs reverse-polarity, TVS, and fusing/PTC
consideration. See Q4.
\item \textbf{Batteries}: no battery on this board (see Section 3).
\item \textbf{RF transmit}: none identified in the frame or brief.
\item \textbf{Mains voltage}: no ? confirmed not applicable (VM is a low-voltage DC
supply rail, not mains-connected).
\end{itemize}
\subsubsection*{Assumptions inferred from the reference frame}
Every item here comes from one frame of the reference video or is our own guess, not a published spec. Research with a source overrides a frame guess (see research/reference-video.md); the comparison is judged against this list.

\begin{itemize}
\item A1 Topology: three half-bridges of N-FETs (Q1-Q6), low-side shunts per phase (SHUNT A/B/C). Frame.
\item A2 Supply: one VM rail with a large electrolytic; 12-24 V nominal, 10-28 V operating. Guess (no rating visible).
\item A3 Power stage rating: 60 V FETs (ours; the frame suggested 40 V class, raised for TVS clamp margin). Guess.
\item A4 Phase current: 10 A RMS continuous, 20 A peak, no heatsink. Guess, from the FET package size.
\item A5 Logic rail: a buck to 3.3 V (L1, BUCK\_SW, V3V3 are visible). Frame; our tree adds 5 V and a gate rail.
\item A6 MCU: a QFP of about 48-64 pins; we chose STM32G4. Frame (package) + guess (family).
\item A7 Gate driver and current-sense amps: some of U1-U10 (SOIC/TSSOP). Frame (packages) + guess (roles).
\item A8 Control I/O: a small JST connector (we made it UART) and a wider connector (we made it Hall + ABZ). Guess.
\item A9 Power/motor connections: TP1-TP5 read as VM, GND, A, B, C wire pads. Guess.
\item A10 User I/O: two tactile switches (SW1/SW2, we made them reset + user), status LEDs. Frame + guess.
\item A11 Stack-up: 4 layers. Frame (inner-plane look), unconfirmed.
\item A12 Size: roughly square, about 60 x 60 mm; mounting holes H1-H10 visible, we use four M3. Frame + guess.
\end{itemize}
\subsubsection*{9. Open questions}
1. What protocol does the small JST-style control connector carry ? UART, CAN, or step/direction + PWM/enable? DEFAULT: UART (simplest, matches a single small connector with few pins) at 3.3V logic levels. 2. What is the Hall/encoder connector's exact pinout and pin count ? 5V, GND, H1-H3 (3 Hall sensors) vs. A/B/Z quadrature encoder? DEFAULT: 5-pin, 5V + GND + H1-H3 (three Hall-effect sensor inputs), matching "sensored" BLDC commutation, the simpler and more common case for a demo board. 3. Confirm the PROVISIONAL current rating: 10 A RMS per phase continuous (no heatsink, 25C ambient), 20 A peak. DEFAULT: accept as stated (orchestrator-provided); flagged here only because Section 8 requires every safety-relevant value to be explicitly confirmed, not assumed. 4. Confirm the PROVISIONAL voltage rating: VM 10-28 V operating (12-24 V nominal), 40 V-rated power stage. DEFAULT: accept as stated (orchestrator-provided); same confirm-not-guess reasoning as Q3. 5. Does the gate-drive rail need its own regulator (independent of VM), or is bootstrap-only sufficient across the full 10-28 V VM range? DEFAULT: bootstrap-only for the high sides (standard for this VM range with a modern 3-phase gate driver IC), with the driver's own internal VCC/VGS rail generated from VM or 3V3 as the driver IC requires ? final choice deferred to P2 part selection, not a P0 blocker. 6. Build quantity and target unit cost? DEFAULT: prototype run, 3-5 boards, no target unit cost (this is a reverse-engineering exercise, not a cost-optimized product). 7. Single-sided or double-sided SMD placement? DEFAULT: single-sided SMD (all fine-pitch parts on top), THT (electrolytic, connectors) also on top ? simplest and cheapest JLCPCB PCBA option, matches "SMD assembly at JLCPCB" in the brief with no double-sided call-out. 8. Exact mounting hole count, pattern, and size? DEFAULT: 4 corner holes, M3, matching the frame's visible corner pads ? remaining edge marks in the frame (H5-H10-style labels) are read as test points/silkscreen refs, not additional mounting holes, pending layout-stage (P5) confirmation.

\subsubsection*{Answers (orchestrator as delegate, unattended run, 2026-09-27)}
1. Control connector: UART (3.3 V logic) on a 4-pin JST-GH: GND, 5V out(fused/limited), TX, RX. Accepted default. 2. Sensor connector: the wider connector is 8-pin: 5V, GND, HA, HB, HC, ENC\_A, ENC\_B, ENC\_Z (Hall and ABZ both supported, 5 V-tolerant inputs with pull-ups). Changed from the 5-pin default so both sensor types work. 3. Current 10 A RMS/phase continuous, 20 A peak: PROVISIONAL; the thermal estimate at P2/P8 may derate it and must say so. 4. VM 10-28 V operating (12-24 V nominal), 40 V FETs: PROVISIONAL. 5. Gate drive: the architect picks it, but the requirement is VGS \textgreater{}= 8 V at VM = 10 V and within the driver's and FETs' absolute max at VM = 28 V + transient; a separate gate-drive rail if the driver cannot take VM directly. 6. Quantity 5, prototype, no unit-cost target (owner: "it needs to be good"). 7. Single-sided SMD on top; THT only for the electrolytic/connectors if unavoidable. 8. Four M3 corner holes. SWD on a 1.27 mm 2x5 (or Tag-Connect footprint) header; the architect chooses.

\subsubsection*{High-voltage clearance scope (owner, 2026-09-28)}
The gate, boot, phase and VM nets keep 0.4 mm (up to 45-50 V) or 0.5 mm (57 V boot/gate) to other copper, with two exceptions, written into the \texttt{.kicad\_dru} as \texttt{hv\_scope\_{*}} rules: 1. Nets of one phase leg (GHx, BSTx, PHASE\_x, GLx, LS\_SRC\_x and the gate-resistor nets of that leg) use the normal 0.13 mm between each other. Tracks sit under solder mask, and IPC-2221B Table 6-1 B4 (coated external conductors) allows 0.13 mm up to 100 V. 2. Inside the gate driver's (U301) courtyard, pad-to-pad and pad-to-track drop to the package's own pitch, floor 0.13 mm: IPC-2221B A6 (external component terminations) allows 0.13 mm for 31-100 V, and the DRV8300 is rated for these pins at this pitch. The same rule covers the VM\_IN ideal-diode controller's (U101, SOT-23-8, 0.65 mm pitch) courtyard: IPC-2221B Table 6-1 A6 (component leads, up to 100 V) applies, VM\_IN is a 60 V-class rail, and U101 is rated for it at this pitch. The routers take one clearance per net, so they keep the full value; the gate-driver escapes are hand-routed.

\section{Architecture}
\subsection*{Decisions of record (state.json)}
{\small
\begin{longtable}{lp{5.2cm}p{7.2cm}}
\toprule
\textbf{Phase} & \textbf{What} & \textbf{Why} \\
\midrule
\endhead
P0 & Open questions 1-8 answered as delegate (UART JST-GH 4p; 8-pin Hall+ABZ; 10 A RMS/20 A pk and 10-28 V PROVISIONAL; qty 5; top-side SMD; 4x M3) & Unattended run; brief says decide and journal; safety values re-confirmed at H1 \\
P1 & Power stage uses 60 V FETs (2-3 mOhm, 5x6), not 40 V & VM TVS clamps \textasciitilde{}45 V at 28 V standoff; 40 V FETs have no margin under clamp. Correctness over the brief's 40 V guess \\
P1 & No on-board fuse (external 20 A slow fuse required, stated in guide); no hot-swap inrush limit; no brake chopper - firmware limits regen on VM sense & Bench prototype; fuse/brake sized to the user's supply; recorded as PROVISIONAL safety answers for H1 \\
P2 & Gate drive: DRV8300D on a 12 V GVDD boosted from 5 V (TPS61040), not DRV8320H's charge pump; overcurrent via CSA -\textgreater{} MCU comparators -\textgreater{} TIM1 break & DRV8320H guarantees only 6.3 V high-side VGS at VM=10 V, below the 8 V requirement \\
P2 & Key parts: NTMFS5C628NL 60 V FETs, LM5164 buck, INA240A1 G=20, STM32G431CBT6, LM74500 RP, SMCJ28A, 3 mOhm 2512 shunts, 4L JLC04162H-7628A 2 oz outer & architect P2; see architecture/blocks.md \\
P2 & Designing under coverage gap: tvs-clamp, boost, ldo, mcu, isns\_a/b/c (no library checklist); buck gap is inrush only & Research cap and iteration spend: 4 of 6 research slots went to the safety-critical blocks (gate drive, half-bridge, shunt sensing, reverse-polarity FET). These are covered by P1 research (power.md, mcu.md, current-sense-amp.md) and will be checked against datasheet extraction at P3; buck inrush is bounded by the buck's own soft-start behind the RP-FET/bulk-cap path \\
P0 & Board named bldc-motor-driver; spec is designed from a general brief; research found nothing about the reference board beyond the frame (research/reference-video.md) & Owner addenda: no reference-product naming; research first, fall back to the frame-derived spec if nothing usable \\
P2 & research task block-half-bridge-inverter-1 closed over 1 unverified draft record(s): half-bridge-inverter-switch-node-damping & --accept-drafts: the slot ships without this knowledge \\
P2 & designing under coverage gap: all 12 P2 slots (power\_in x2, regulators x3, gate\_driver, power\_stage, sensing, mcu, isns\_a/b/c) & maturity floor is owner-approved; 19 verified research records (4 tasks) fold in only as provisional until promoted, switch-node-damping record refuted (10-48 V envelope unsupported, source is 5-21 V) and closed as draft; regulators/mcu/isns covered by P1 research md + P3 datasheet extraction; unattended run \\
P3 & designing under coverage gap at P3 exit: 9 topology/interface slots (tvs-clamp, buck, boost, ldo, half-bridge-inverter, mcu, isns\_a/b/c) have no checklist; parts 8 covered/3 provisional & No new part-level gaps vs P2; hwde has no checklists for these topologies; the verified research records + datasheet extractions (8 ICs incl. STM32G431) carry the layout rules. Unattended run. \\
P8 & check\_diffpair (ISNS\_A/B/C): exit 1, 6 errors (skew 14.69/7.90/15.14 mm, N uncoupled 17.35/10.56/17.80 mm vs 3.0 mm limit) + 7 warnings (open\_trunk x6: terminals are net ties; B via asymmetry). check\_diffpair has no waiver input; the same 6 findings are waived in reports/verify-waivers.json and the verify gate records PASS. & Owner-approved: INA240 IN+ (pin 8) and IN- (pin 1) are on opposite pad rows 5.4 mm apart and the shunt Kelvin taps 3.46 mm apart, so \textasciitilde{}9 mm uncoupled is set by the packages; INA240 has its own PWM rejection and the signal is low-frequency sense. A/C N legs loop round C511/C513; a B.Cu underpass would cut them to \textasciitilde{}10.5 mm but add two vias and still fail, so left. \\
P9 & Fab package exported to fab/ (gerbers+drill zip, pos, BOM.csv, CPL.csv, prebuy.csv); BOM incomplete: 53 placed parts (sensing U501-503/C5xx/R5xx, io C70x/D70x/F70x/J70x/R70x) carry no LCSC number on the board and parts.json has no per-ref map. Not ordered. & P9 export per brief; LCSC back-annotation needs schematic LCSC fields, left for the sourcing step rather than guessed. \\
P9 & Fab package re-exported from the final board (pcb sexpr\_no\_uuid:b620b16a..., after the 08:12-08:15 erc/drc\_routed/place/verify/dfm PASS runs): gerbers, drill, pos, BOM.csv, BOM-full.csv, CPL.csv, prebuy.csv. BOM complete: every placed part carries an LCSC number (bom\_cpl pass, 0 missing). Supersedes the 07:35 P9 entry that said 53 parts had none. Not ordered. & The 53 sensing/io parts had LCSC numbers only on their embedded lib symbols ('LCSC Part'); copying them to the instances and footprints as 'LCSC' changed only hidden fields, so the gerber design hash is unchanged (97bbca5e...) and the copper is the one that passed DRC. \\
\bottomrule
\end{longtable}
}
\subsection*{Sheet plan}
\subsection*{hierarchical sheet plan (P2)}
Net naming (canonical; constraints.json uses these):

\begin{itemize}
\item Power nets are global and bare: \texttt{VM}, \texttt{+12V}, \texttt{+5V}, \texttt{+3V3}, \texttt{GND}.
\item Cross-sheet interface nets are labelled on the ROOT sheet, so they come out as \texttt{/NAME}.
\item Sheet-local nets come out as \texttt{/\textless{}sheet\textgreater{}/NAME}.
\item One GND net only (no AGND/PGND split); separation is done by placement and Kelvin routing.
\end{itemize}
Refdes: every sheet owns a hundred-block (R101..R199, C101.., U101.., \#PWR0101..); pwr\_base = the sheet base. Root-sheet parts: H1-H4 (M3 corner holes), fiducials FID1-FID3 if P5 wants them.

\begin{quote}\small\ttfamily\raggedright
\textbar{} Sheet \textbar{} Base \textbar{} Blocks \textbar{} Interface nets (hier pins) \textbar{}\\
\textbar{}---\textbar{}---\textbar{}---\textbar{}---\textbar{}\\
\textbar{} power\_in \textbar{} 100 \textbar{} VM/GND pads TP101 (VM), TP102 (GND); Q101 RP FET; U101 LM74500; D101 SMCJ28A; C101, C104 470uF/63V THT; C102-C103 MLCC \textbar{} out: `VM`, `GND` \textbar{}\\
\textbar{} regulators \textbar{} 200 \textbar{} U201 LM5164 + L201; U202 TPS61040 + L202 + D201; U203 TLV75733; FB201 VDDA ferrite \textbar{} in: `VM`; out: `+5V`, `+12V`, `+3V3`, `/VDDA` \textbar{}\\
\textbar{} gate\_driver \textbar{} 300 \textbar{} U301 DRV8300D; C301-C303 boot caps; GVDD caps; R301 DT, MODE strap \textbar{} in: `+12V`, `/INHA` `/INLA` `/INHB` `/INLB` `/INHC` `/INLC`; out: `/GHA` `/GLA` `/GHB` `/GLB` `/GHC` `/GLC`; bidir: `/PHASE\_A` `/PHASE\_B` `/PHASE\_C` (SHx) \textbar{}\\
\textbar{} power\_stage \textbar{} 400 \textbar{} leg A: Q401 (HS), Q402 (LS), R401 shunt, C401-C403 10uF 1210, C404 100nF; leg B: Q403, Q404, R402, C405-C408; leg C: Q405, Q406, R403, C409-C412; gate R R411-R416; gate pull-downs R421-R426; DNP snubbers R431-R436 / C431-C436; phase pads TP401-TP403 \textbar{} in: `VM`, `GND`, `/GHx`, `/GLx`; out: `/PHASE\_A/B/C`, `/ISNS\_A\_P` `/ISNS\_A\_N` (and B, C) \textbar{}\\
\textbar{} sensing \textbar{} 500 \textbar{} U501-U503 INA240A1 + input RC; phase dividers R5xx for `/VSENSE\_A/B/C`; bus divider for `/VBUS\_SENSE` \textbar{} in: `/ISNS\_x\_P/N`, `/PHASE\_x`, `VM`, `+3V3`; out: `/ISENSE\_A/B/C`, `/VSENSE\_A/B/C`, `/VBUS\_SENSE` \textbar{}\\
\textbar{} mcu \textbar{} 600 \textbar{} U601 STM32G431CBT6 + decoupling; J601 SWD 2x5 1.27; SW601 reset, SW602 user; D601-D603 LEDs (power, status, fault) \textbar{} in: `+3V3`, `/VDDA`, `/ISENSE\_x`, `/VSENSE\_x`, `/VBUS\_SENSE`, `/HALL\_A/B/C`, `/ENC\_A/B/Z`, `/UART\_RX`; out: `/INHx`, `/INLx`, `/UART\_TX` \textbar{}\\
\textbar{} io \textbar{} 700 \textbar{} J701 JST-GH 4 control; J702 8-pin sensor; F701/F702 0.2 A PTC; ESD + series R + pull-ups + RC filters \textbar{} in: `+5V`, `+3V3`; bidir: `/UART\_TX`, `/UART\_RX`, `/HALL\_A/B/C`, `/ENC\_A/B/Z` \textbar{}
\end{quote}
Sheet-local nets that constraints.json names: \texttt{/power\_in/VM\_IN} (pad to RP FET), \texttt{/power\_stage/LS\_SRC\_A/B/C} (LS source = shunt top), \texttt{/gate\_driver/BSTA/B/C}, \texttt{/regulators/BUCK\_SW}, \texttt{/regulators/BOOST\_SW}. Connector-side 5 V after the PTCs: \texttt{/io/CTL\_5V}, \texttt{/io/SEN\_5V}.

MCU pin notes for P3/P4: TIM1 CH1-3 + CH1N-3N for INH/INL (PA8/PA9/PA10, PB13/PB14/PB15); ISENSE\_x on pins that are both ADC and COMP+ inputs (COMP -\textgreater{} TIM1 BKIN for overcurrent); Halls on a timer's CH1-3 in XOR/hall mode; encoder on another timer's CH1/CH2 + ETR/GPIO for Z; UART on USART1 or USART2; PB8/BOOT0 kept for BOOT0 (pull-down), not used as an input.

Silkscreen: "PCB-0018 REV A" on F.Silk, plus pad labels VM, GND, A, B, C beside the power pads.

Full architecture narrative (not inlined; see the artifact index): \texttt{architecture/blocks.md}, \texttt{architecture/power\_tree.md}, \texttt{architecture/stackup.md}.
\section{Schematic}
Gate \texttt{erc} (phase P4): \textbf{pass} after 10 attempt(s); last 2026-09-28T08:13:54, failing 0 of 0.
The full schematic PDF follows.
\includepdf[pages=-]{reports/schematic.pdf}
\section{Layout}
Gate \texttt{place} (phase P6): \textbf{pass} after 10 attempt(s); last 2026-09-28T08:14:07, failing 0 of 14.

Gate \texttt{drc\_routed} (phase P7): \textbf{pass} after 5 attempt(s); last 2026-09-28T08:14:02, failing 0 of 0.
\begin{center}
\includegraphics[width=0.72\textwidth]{reports/PCB-0018-A_bldc-motor-driver_top.png}\\
{\small\texttt{reports/PCB-0018-A\_bldc-motor-driver\_top.png}}
\end{center}
\begin{center}
\includegraphics[width=0.72\textwidth]{reports/PCB-0018-A_bldc-motor-driver_bottom.png}\\
{\small\texttt{reports/PCB-0018-A\_bldc-motor-driver\_bottom.png}}
\end{center}
\section{Routing Renders}
The finished board's renders are under Layout; this section shows how it got there.
 The run kept no routing snapshots.
\section{Design Decisions}
Every decision the run recorded in state.json (13), oldest first, with the reason given at the time.
{\small
\begin{longtable}{lp{1.6cm}p{6.2cm}p{6.2cm}}
\toprule
\textbf{Ph} & \textbf{When} & \textbf{Decision} & \textbf{Why} \\
\midrule
\endhead
P0 & 09-27 14:04 & Open questions 1-8 answered as delegate (UART JST-GH 4p; 8-pin Hall+ABZ; 10 A RMS/20 A pk and 10-28 V PROVISIONAL; qty 5; top-side SMD; 4x M3) & Unattended run; brief says decide and journal; safety values re-confirmed at H1 \\
P1 & 09-27 14:08 & Power stage uses 60 V FETs (2-3 mOhm, 5x6), not 40 V & VM TVS clamps \textasciitilde{}45 V at 28 V standoff; 40 V FETs have no margin under clamp. Correctness over the brief's 40 V guess \\
P1 & 09-27 14:08 & No on-board fuse (external 20 A slow fuse required, stated in guide); no hot-swap inrush limit; no brake chopper - firmware limits regen on VM sense & Bench prototype; fuse/brake sized to the user's supply; recorded as PROVISIONAL safety answers for H1 \\
P2 & 09-27 14:26 & Gate drive: DRV8300D on a 12 V GVDD boosted from 5 V (TPS61040), not DRV8320H's charge pump; overcurrent via CSA -\textgreater{} MCU comparators -\textgreater{} TIM1 break & DRV8320H guarantees only 6.3 V high-side VGS at VM=10 V, below the 8 V requirement \\
P2 & 09-27 14:26 & Key parts: NTMFS5C628NL 60 V FETs, LM5164 buck, INA240A1 G=20, STM32G431CBT6, LM74500 RP, SMCJ28A, 3 mOhm 2512 shunts, 4L JLC04162H-7628A 2 oz outer & architect P2; see architecture/blocks.md \\
P2 & 09-27 14:32 & Designing under coverage gap: tvs-clamp, boost, ldo, mcu, isns\_a/b/c (no library checklist); buck gap is inrush only & Research cap and iteration spend: 4 of 6 research slots went to the safety-critical blocks (gate drive, half-bridge, shunt sensing, reverse-polarity FET). These are covered by P1 research (power.md, mcu.md, current-sense-amp.md) and will be checked against datasheet extraction at P3; buck inrush is bounded by the buck's own soft-start behind the RP-FET/bulk-cap path \\
P0 & 09-27 14:32 & Board named bldc-motor-driver; spec is designed from a general brief; research found nothing about the reference board beyond the frame (research/reference-video.md) & Owner addenda: no reference-product naming; research first, fall back to the frame-derived spec if nothing usable \\
P2 & 09-27 14:55 & research task block-half-bridge-inverter-1 closed over 1 unverified draft record(s): half-bridge-inverter-switch-node-damping & --accept-drafts: the slot ships without this knowledge \\
P2 & 09-27 14:56 & designing under coverage gap: all 12 P2 slots (power\_in x2, regulators x3, gate\_driver, power\_stage, sensing, mcu, isns\_a/b/c) & maturity floor is owner-approved; 19 verified research records (4 tasks) fold in only as provisional until promoted, switch-node-damping record refuted (10-48 V envelope unsupported, source is 5-21 V) and closed as draft; regulators/mcu/isns covered by P1 research md + P3 datasheet extraction; unattended run \\
P3 & 09-27 15:44 & designing under coverage gap at P3 exit: 9 topology/interface slots (tvs-clamp, buck, boost, ldo, half-bridge-inverter, mcu, isns\_a/b/c) have no checklist; parts 8 covered/3 provisional & No new part-level gaps vs P2; hwde has no checklists for these topologies; the verified research records + datasheet extractions (8 ICs incl. STM32G431) carry the layout rules. Unattended run. \\
P8 & 09-28 07:26 & check\_diffpair (ISNS\_A/B/C): exit 1, 6 errors (skew 14.69/7.90/15.14 mm, N uncoupled 17.35/10.56/17.80 mm vs 3.0 mm limit) + 7 warnings (open\_trunk x6: terminals are net ties; B via asymmetry). check\_diffpair has no waiver input; the same 6 findings are waived in reports/verify-waivers.json and the verify gate records PASS. & Owner-approved: INA240 IN+ (pin 8) and IN- (pin 1) are on opposite pad rows 5.4 mm apart and the shunt Kelvin taps 3.46 mm apart, so \textasciitilde{}9 mm uncoupled is set by the packages; INA240 has its own PWM rejection and the signal is low-frequency sense. A/C N legs loop round C511/C513; a B.Cu underpass would cut them to \textasciitilde{}10.5 mm but add two vias and still fail, so left. \\
P9 & 09-28 07:35 & Fab package exported to fab/ (gerbers+drill zip, pos, BOM.csv, CPL.csv, prebuy.csv); BOM incomplete: 53 placed parts (sensing U501-503/C5xx/R5xx, io C70x/D70x/F70x/J70x/R70x) carry no LCSC number on the board and parts.json has no per-ref map. Not ordered. & P9 export per brief; LCSC back-annotation needs schematic LCSC fields, left for the sourcing step rather than guessed. \\
P9 & 09-28 08:38 & Fab package re-exported from the final board (pcb sexpr\_no\_uuid:b620b16a..., after the 08:12-08:15 erc/drc\_routed/place/verify/dfm PASS runs): gerbers, drill, pos, BOM.csv, BOM-full.csv, CPL.csv, prebuy.csv. BOM complete: every placed part carries an LCSC number (bom\_cpl pass, 0 missing). Supersedes the 07:35 P9 entry that said 53 parts had none. Not ordered. & The 53 sensing/io parts had LCSC numbers only on their embedded lib symbols ('LCSC Part'); copying them to the instances and footprints as 'LCSC' changed only hidden fields, so the gerber design hash is unchanged (97bbca5e...) and the copper is the one that passed DRC. \\
\bottomrule
\end{longtable}
}
\section{Design History}
\subsection*{Phase timeline}
{\small
\begin{longtable}{llllr}
\toprule
\textbf{Phase} & \textbf{Name} & \textbf{First decision} & \textbf{Last decision} & \textbf{Decisions} \\
\midrule
\endhead
P0 & Intake & 2026-09-27T14:04 & 2026-09-27T14:32 & 2 \\
P1 & Research & 2026-09-27T14:08 & 2026-09-27T14:08 & 2 \\
P2 & Architecture & 2026-09-27T14:26 & 2026-09-27T14:56 & 5 \\
P3 & Parts and library & 2026-09-27T15:44 & 2026-09-27T15:44 & 1 \\
P8 & Verification & 2026-09-28T07:26 & 2026-09-28T07:26 & 1 \\
P9 & DFM & 2026-09-28T07:35 & 2026-09-28T08:38 & 2 \\
\bottomrule
\end{longtable}
}
\subsection*{Human checkpoints}
{\small
\begin{longtable}{lllp{9cm}}
\toprule
\textbf{H} & \textbf{Status} & \textbf{When} & \textbf{Note} \\
\midrule
\endhead
1 & skipped & 2026-09-27T14:56 & Unattended run (owner: decide for yourself): architecture taken as PROVISIONAL delegate answer, not owner approval. 10-28 V, 10 A RMS/20 A pk, 60 V FETs, DRV8300D bootstrap, 3x Kelvin shunts + INA240, STM32G4, LM74500 RP NFET, buck 5V/LDO 3V3/12V boost, 4L. All P2 slots under coverage gap. \\
\bottomrule
\end{longtable}
}
\subsection*{Where the run went back}
\paragraph*{P2 back to P0, 2026-09-27T14:32}
\begin{itemize}
\item Board named bldc-motor-driver; spec is designed from a general brief; research found nothing about the reference board beyond the frame (research/reference-video.md) \emph{Why:} Owner addenda: no reference-product naming; research first, fall back to the frame-derived spec if nothing usable
\end{itemize}
\subsection*{Component choices and changes}
{\small
\begin{longtable}{lp{1.6cm}p{6.2cm}p{6.2cm}}
\toprule
\textbf{Ph} & \textbf{When} & \textbf{Decision} & \textbf{Why} \\
\midrule
\endhead
P1 & 09-27 14:08 & No on-board fuse (external 20 A slow fuse required, stated in guide); no hot-swap inrush limit; no brake chopper - firmware limits regen on VM sense & Bench prototype; fuse/brake sized to the user's supply; recorded as PROVISIONAL safety answers for H1 \\
P2 & 09-27 14:26 & Gate drive: DRV8300D on a 12 V GVDD boosted from 5 V (TPS61040), not DRV8320H's charge pump; overcurrent via CSA -\textgreater{} MCU comparators -\textgreater{} TIM1 break & DRV8320H guarantees only 6.3 V high-side VGS at VM=10 V, below the 8 V requirement \\
P0 & 09-27 14:32 & Board named bldc-motor-driver; spec is designed from a general brief; research found nothing about the reference board beyond the frame (research/reference-video.md) & Owner addenda: no reference-product naming; research first, fall back to the frame-derived spec if nothing usable \\
\bottomrule
\end{longtable}
}
\subsection*{Git history}
37 commits touched the workspace, 2026-09-27T14:33 to 2026-09-28T08:30. By hour, with the files each area had changed:
{\small
\begin{longtable}{lrp{10cm}}
\toprule
\textbf{Hour} & \textbf{Commits} & \textbf{Areas (files)} \\
\midrule
\endhead
2026-09-27 14:00 & 3 & research 76, architecture 5, log 4, (root) 3, brief 2 \\
2026-09-27 15:00 & 7 & lib 101, parts 31, research 23, reports 8, log 6, (root) 5, kicad 3, learnings 1, architecture 1 \\
2026-09-27 16:00 & 3 & kicad 24, log 3, reports 2, (root) 2, architecture 1, lib 1, parts 1 \\
2026-09-27 17:00 & 1 & kicad 18, log 5, architecture 3, (root) 3, anneal 2, parts 1 \\
2026-09-27 18:00 & 1 & anneal 1, kicad 1, log 1 \\
2026-09-27 19:00 & 1 & kicad 2, log 2, (root) 2 \\
2026-09-27 20:00 & 3 & log 6, (root) 3, kicad 3, architecture 1 \\
2026-09-27 23:00 & 1 & log 2, kicad 1 \\
2026-09-28 00:00 & 2 & log 2, kicad 2, (root) 1 \\
2026-09-28 01:00 & 1 & log 2, kicad 1 \\
2026-09-28 02:00 & 1 & kicad 1 \\
2026-09-28 03:00 & 2 & kicad 3, routing 1 \\
2026-09-28 04:00 & 2 & kicad 2, routing 2 \\
2026-09-28 05:00 & 3 & kicad 3, routing 3 \\
2026-09-28 06:00 & 3 & kicad 9, (root) 7, routing 2, reports 2 \\
2026-09-28 07:00 & 2 & fab 21, reports 4, (root) 4, kicad 2, routing 2 \\
2026-09-28 08:00 & 1 & fab 19, reports 14, kicad 4, (root) 3 \\
\bottomrule
\end{longtable}
}
\subsection*{Comparison}
\subsection*{PCB-0018-A bldc-motor-driver vs the reference board}
The reference is a 3-phase motor-driver board seen in one frame of a published demo video (\texttt{brief/reference-video-frame.png}). Nobody published its spec, schematic or parts list (research/reference-video.md), so the left column is only what the frame shows. Our spec was written from that frame plus our own guesses; every guess is listed as A1-A12 in requirements.md.

\begin{quote}\small\ttfamily\raggedright
\textbar{} \textbar{} Reference board (from the frame) \textbar{} hwde, PCB-0018 rev A \textbar{}\\
\textbar{}---\textbar{}---\textbar{}---\textbar{}\\
\textbar{} Function \textbar{} 3-phase motor driver \textbar{} 3-phase BLDC/PMSM driver, sensored or sensorless FOC \textbar{}\\
\textbar{} Power stage \textbar{} Q1-Q6 in dual-FET-looking packages, three half-bridges \textbar{} 6x NTMFS5C628NL 60 V, 5x6 mm, one FET per package \textbar{}\\
\textbar{} Current sensing \textbar{} low-side shunts SHUNT A/B/C \textbar{} 3 mOhm 2512 low-side shunts, Kelvin taps via net ties, 3x INA240A1 \textbar{}\\
\textbar{} Supply \textbar{} VM rail, one large electrolytic, no rating visible \textbar{} VM 10-28 V (12-24 V nominal), 2x 470 uF / 63 V, SMCJ28A TVS, reverse-polarity FET \textbar{}\\
\textbar{} Logic rail \textbar{} buck to 3.3 V (L1, BUCK\_SW, V3V3) \textbar{} LM5164 buck VM-\textgreater{}5 V, TLV75733 LDO 5-\textgreater{}3.3 V, TPS61040 boost 5-\textgreater{}12 V for gate drive \textbar{}\\
\textbar{} Gate driver \textbar{} one of U1-U10, not identifiable \textbar{} DRV8300D, 6-PWM, bootstrap high sides \textbar{}\\
\textbar{} MCU \textbar{} QFP, about 48-64 pins \textbar{} STM32G431CBT6 (LQFP-48) \textbar{}\\
\textbar{} Control I/O \textbar{} small JST connector, wider FFC/header \textbar{} JST-GH 4 (UART + 5 V), 8-pin Hall + ABZ encoder, SWD 2x5 1.27 mm \textbar{}\\
\textbar{} User I/O \textbar{} SW1, SW2 \textbar{} reset + user switch, 3 LEDs \textbar{}\\
\textbar{} Test points / wire pads \textbar{} TP1-TP5 \textbar{} TP1-TP5 as VM, GND, A, B, C wire pads \textbar{}\\
\textbar{} Mounting \textbar{} H1-H10 \textbar{} 4x M3 corner holes \textbar{}\\
\textbar{} Layers \textbar{} looks like 4 \textbar{} 4 (JLC04162H-7628A, 2 oz outer) \textbar{}\\
\textbar{} Size \textbar{} roughly square, about 60 x 60 mm (guess) \textbar{} 78 x 78 mm \textbar{}\\
\textbar{} Footprints \textbar{} not countable from the frame \textbar{} 175: 148 machine-placed, 12 do-not-place, 15 board features \textbar{}\\
\textbar{} ERC / DRC \textbar{} not reported \textbar{} ERC and DRC 0 violations, 0 unconnected; see below for what was waived \textbar{}\\
\textbar{} Fabricated / measured \textbar{} not reported \textbar{} no (fab package only, nothing ordered) \textbar{}
\end{quote}
\subsubsection*{What the checks say about ours}
ERC, DRC, placement, the design checks and DFM all pass, and the fab package has gerbers, drill, a BOM with an LCSC number on every placed part (45 lines) and a placement file. The design checks pass with 390 waivers, each with its reason in \texttt{routing/rules-added.md}, and you should read them before trusting the board:

\begin{itemize}
\item \textbf{Current (359).} Most are stitching vias and Kelvin or gate-return taps that the check counts as
current paths; the real phase and supply necks were widened or given more vias.
\item \textbf{Creepage (24).} They sit inside the gate driver's and the ideal-diode controller's own pin pitch
(IPC-2221B A6, leads up to 100 V) or between nets of the same phase leg.
\item \textbf{Current-sense pairs (6).} The INA240's two inputs sit 5.4 mm apart on opposite pad rows, so
each pair runs at least 9 mm uncoupled against a 3 mm target. The packages set that length, not
the routing, and the traces run side by side everywhere else.
\item \textbf{Power distribution (1).} The check flags the ground rail itself as undecoupled.
\end{itemize}
Some narrow power stubs were turned into solid copper zones of the same shape to clear the track-width rule, so those necks are no wider than before. Nothing here has been built or measured.

\subsubsection*{Where we deviate, and why}
\begin{itemize}
\item \textbf{Single FETs, not dual packages.} A dual-FET package puts both switches of a leg in one thermal
mass. At our 10 A RMS guess, two 5x6 mm single FETs spread the heat over twice the copper, and
each leg's high-frequency capacitors sit right across its drain and shunt.
\item \textbf{Three rails, not one.} The frame only shows a 3.3 V buck. A 3-phase driver also needs a gate
rail of about 12 V that holds up when VM sags to 10 V, so we added a 5 V buck and a 12 V boost.
\item \textbf{Bigger board.} Three legs of power copper plus the reverse-polarity FET, a 470 uF can pair and the 100 V buck don't fit in
60 x 60 mm without crowding the gate loops. We put gate-loop and power-loop area first and let size
come out at 78 x 78 mm.
\item \textbf{Four mounting holes, not ten.} Most of the frame's H labels look like fiducials or edge marks
at that resolution; we read them as optional.
\end{itemize}
\section{How the Design Process Went}
Each phase the run passed through, top to bottom; a red box is a point where it went back to an earlier phase, with the first decision that sent it there. The current phase has the red outline.
\begin{center}
\includegraphics[width=0.8\textwidth]{reports/design_full/flow.pdf}\\
{\small\texttt{reports/design\_full/flow.pdf}}
\end{center}
\section{Verification}
Gate \texttt{verify} (phase P8): \textbf{pass} after 11 attempt(s); last 2026-09-28T08:15:06, failing 0 of 450.
\emph{reports/verify\_all.json not available.}
\section{DFM and Fabrication}
Gate \texttt{dfm} (phase P9): \textbf{pass} after 3 attempt(s); last 2026-09-28T08:15:25, failing 0 of 4.
\subsection*{Fabrication outputs}
Exported layers: F.Cu, In1.Cu, In2.Cu, B.Cu, F.Silkscreen, B.Silkscreen, F.Mask, B.Mask, F.Paste, B.Paste, Edge.Cuts.
\subsection*{Bill of materials}
{\small
\begin{longtable}{p{4.4cm}p{2.8cm}p{4.6cm}l}
\toprule
\textbf{Comment} & \textbf{Designator} & \textbf{Footprint} & \textbf{LCSC} \\
\midrule
\endhead
470uF 63V & C101,C104 & CAP-TH\_BD13.0-P5.00-D1.2-FD & C106675 \\
100nF 50V X7R & C102,C103,C203,C209,C211,C213,C511,C512,C513,C524,C601,C602,C603,C608,C609,C610 & C0603 & C14663 \\
2.2nF 50V X7R & C201,C210,C501,C502,C503,C521,C522,C523,C701,C702,C703,C704,C705,C706 & C0603 & C1604 \\
22uF 25V X5R & C202 & C0805 & C45783 \\
10uF 50V X7R & C204,C304,C401,C402,C403,C405,C406,C407,C409,C410,C411 & C1210 & C2918502 \\
10uF 25V X5R & C205 & C0805 & C15850 \\
1uF 50V X7R & C206,C207,C208,C301,C302,C303,C305,C605,C607 & C0805 & C28323 \\
1nF 50V X7R & C212 & C0603 & C1588 \\
100nF 50V X7R & C404,C408,C412 & C1210 & C76567 \\
4.7uF 25V X5R & C604 & C0805 & C1779 \\
10nF 50V X7R & C606 & C0603 & C57112 \\
SMCJ28A & D101 & SMC\_L6.9-W5.8-LS7.9-RD & C2943858 \\
MBR0520L & D201 & SOD-123\_L2.7-W1.6-LS3.7-RD-1 & C5563735 \\
LED green (power) & D601 & LED-SMD\_L1.6-W0.8-R-RD & C12624 \\
LED yellow (status) & D602 & LED-SMD\_L1.6-W0.8-R-RD & C2287 \\
LED red (fault) & D603 & LED-SMD\_L1.6-W0.8-R-RD & C2286 \\
PESD5V0S1BA-ES & D701,D702,D703,D704,D705,D709,D710,D711 & SOD-323\_L1.7-W1.3-LS2.6-BI & C5180216 \\
PTC 0.2A hold 30V & F701,F702 & F1206 & C69680 \\
600ohm@100MHz ferrite & FB201 & R0603 & C1002 \\
SWD 2x5 1.27mm & J601 & CONN-SMD\_10P-P1.27\_XFCN\_PZ127VS-12-10-H10-A38 & C2935458 \\
JST-GH 4p SMD & J701 & CONN-SMD\_4P-P1.25\_SM04B-GHS-TB-LF-SN & C189895 \\
JST-GH 8p SMD & J702 & CONN-SMD\_SH-SM08B-GHS-TB & C42376901 \\
22uH shielded & L201 & IND-SMD\_L6.0-W6.0 & C2849539 \\
22uH shielded & L202 & IND-SMD\_L4.0-W4.0\_LQH44PN2R2MP0L & C135264 \\
NTMFS5C628NL & Q101,Q401,Q402,Q403,Q404,Q405,Q406 & DFN-5\_L4.9-W5.9-P1.27-LS6.2-BL & C145537 \\
31.6k 1\% & R201 & R0603 & C25967 \\
33.2k 1\% & R202 & R0603 & C23003 \\
10k & R203,R205,R302,R421,R422,R423,R424,R425,R426,R511,R513,R515,R517,R601 & R0603 & C25804 \\
86.6k 1\% & R204 & R0603 & C23121 \\
100k & R206 & R0603 & C25803 \\
40.2k 1\% & R301 & R0603 & C12447 \\
3mR 2512 3W 1\% & R401,R402,R403 & R2512 & C5375459 \\
10R & R411,R412,R413,R414,R415,R416 & R0603 & C22859 \\
100 ohm & R501,R502,R503,R701,R702,R703,R704,R705,R709,R710,R711 & R0603 & C22775 \\
150k & R510,R512,R514,R516 & R0603 & C22807 \\
1k 1\% & R602,R603,R604 & R0603 & C21190 \\
4.7k & R706,R707,R708 & R0603 & C23162 \\
Reset switch SMD & SW601,SW602 & SW-SMD\_L3.9-W3.0-P4.45 & C720477 \\
LM74500-Q1 & U101 & SOT-23-8\_L3.0-W1.7-P0.65-LS2.8-BR & C5219210 \\
LM5164 & U201 & SO-8\_L4.9-W3.9-P1.27-LS6.0-BL-EP & C477928 \\
TPS61040 & U202 & SOT-23-5\_L2.9-W1.6-P0.95-LS2.8-BL & C7722 \\
TLV75733 & U203 & SOT-23-5\_L2.9-W1.6-P0.95-LS2.8-BL & C485517 \\
DRV8300DRGER & U301 & VQFN-24\_L4.0-W4.0-P0.50-TL-EP2.5 & C3655801 \\
INA240A1 & U501,U502,U503 & SOIC-8\_L5.0-W4.0-P1.27-LS6.0-BL & C2060769 \\
STM32G431CBT6 & U601 & LQFP-48\_L7.0-W7.0-P0.50-LS9.0-BL & C529355 \\
\bottomrule
\end{longtable}
}
15 CPL rotation correction(s) applied; no missing LCSC numbers.
\emph{fab/order.json not available.}
\section{Run Record (Appendix)}
\subsection*{Phase digests}
\paragraph*{P0}
\begin{quote}\small\ttfamily\raggedright
\# P0 digest\\
- requirements.md by requirements-analyst (sonnet); spec inferred from the reference video frame, no mode token.\\
- 8 open questions answered by the orchestrator as delegate; 10 A RMS / 20 A pk and VM 10-28 V are PROVISIONAL.\\
- 4 layers; size is an output, \textasciitilde{}60 x 60 mm soft target, roughly square; 4x M3.\\
- Correctness order: gate drive, Kelvin shunts, power loop, thermal, then size.
\end{quote}
\paragraph*{P1}
\begin{quote}\small\ttfamily\raggedright
\# P1 digest\\
- power (opus): RP high-side NFET + LM74500-class; TVS 28 V standoff (\textasciitilde{}45 V clamp); buck VM-\textgreater{}5V, LDO 3V3, boost 12 V GVDD; 9x10uF 1210 + \textgreater{}=470uF bulk; \textasciitilde{}3.5 W board loss.\\
- refdesign power stage: 10R gates, 1uF boot, Kelvin shunts 3-5 mOhm, 1:11 dividers, per-leg ceramics, TIDA-010251 thermal \textasciitilde{}49 C at 10 A RMS.\\
- scouts (ran before power.md existed): DRV8320HRTVR driver (6-60 V, own charge pump -\textgreater{} the 12 V boost may be unnecessary),\\
  TDM3436 40 V FET, INA181A2 CSA, STM32G431CBT6, TPS54331 buck.\\
- CONFLICTS for P2 architect: FET must be 60 V (decision); buck/any IC on VM must take \textgreater{}= 45 V clamp (TPS54331 is 28 V - reject);\\
  DRV8320 charge pump vs power.md's 12 V boost - pick one.\\
- Decisions: 60 V FETs; external fuse, no hot-swap, no brake chopper (PROVISIONAL).
\end{quote}
\paragraph*{P2}
\begin{quote}\small\ttfamily\raggedright
\# P2 digest\\
- architect (opus) wrote architecture/\{blocks,power\_tree,stackup,sheets\}.md + constraints.json (lint 0/0).\\
- DRV8300D + 12 V boost GVDD (DRV8320H rejected: 6.3 V VGS at VM=10 V). OC via INA240 -\textgreater{} G4 COMP -\textgreater{} TIM1 BKIN.\\
- 60 V NTMFS5C628NL x6, LM5164 100 V buck, LM74500 RP (23 in stock - recheck at P3), SMCJ28A, 1:16 dividers.\\
- Riskiest: 10 A RMS no heatsink (\textasciitilde{}3.4 W); P8 check\_thermal confirms or derates.\\
- Coverage P2: 12/12 slots gap -\textgreater{} research not yet run.
\end{quote}
\paragraph*{P3}
\emph{(no digest recorded)}
\paragraph*{P4}
\begin{quote}\small\ttfamily\raggedright
\# P4 digest - schematic\\
- 7 sheets generated from kicad/gen/{*}.py (root.py stitches); ERC gate PASS 0/0, netlist\_audit PASS.\\
- schematic-reviewer: 2 errors (LM5164 ripple injection, TP101/102 pad size) fixed; dispositions in reports/review-schematic.md.\\
- Closed OPEN items: C212 1 nF CFF on TPS61040 (datasheet eq. 6); L202 rated 1.05 A vs 0.47 A peak; C104 second 470 uF bulk cap (bus ripple), derating note in requirements.md; NT401-406 in leg groups.\\
- Kelvin: net ties NT401-406 must sit on the shunt pads' inner edge at P6.\\
- DNP: RC snubbers C431-436/R431-436 (placeholders 3.3R/2.2nF).
\end{quote}
\paragraph*{P5}
\begin{quote}\small\ttfamily\raggedright
\# P5 digest - board setup\\
- board\_init PASS: 4L JLC04162H-7628A (2 oz outer), 4layer\_2oz floors, parity 0, setup 0, unconnected 365 (P7's).\\
- Outline auto (shelf pack); target \textasciitilde{}60-70 mm square is set at P6 with board\_edit --outline.\\
- rules\_gen PASS: 43 rules; classes Pwr\_4p2566mm (VM\_IN/VM/GND/phases 13 A), Pwr\_2p9666mm, Pwr\_0p225mm, Diff90 (ISNS pairs); 0.4 mm HV clearance on phase nets.\\
- Host note: stock libs passed via --fp-lib \textasciitilde{}/.local/kicad10/usr/share/kicad/footprints; project fp-lib-table names MountingHole/NetTie/TestPoint via \$\{KICAD10\_FOOTPRINT\_DIR\}.\\
- OPEN for P6/P7: silkscreen "PCB-0018 REV A" gr\_text once the outline is final; ISNS diff pairs at 0.56 mm are wide for the CSA pins - neck down at the IC.
\end{quote}
\paragraph*{P6}
\emph{(no digest recorded)}
\paragraph*{P7}
\emph{(no digest recorded)}
\paragraph*{P8}
\emph{(no digest recorded)}
\paragraph*{P9}
\emph{(no digest recorded)}
\subsection*{Gate attempt history}
{\small
\begin{longtable}{llll}
\toprule
\textbf{Timestamp} & \textbf{Gate} & \textbf{Status} & \textbf{Failing/Total} \\
\midrule
\endhead
2026-09-27T16:09:34 & erc & fail & 9/9 \\
2026-09-27T16:11:16 & erc & pass & 0/0 \\
2026-09-27T16:15:07 & erc & pass & 0/0 \\
2026-09-27T16:25:32 & erc & pass & 0/0 \\
2026-09-27T16:26:29 & erc & pass & 0/0 \\
2026-09-27T16:28:35 & erc & pass & 0/0 \\
2026-09-27T16:35:46 & erc & pass & 0/0 \\
2026-09-27T16:40:22 & erc & pass & 0/0 \\
2026-09-27T17:19:59 & place & fail & 24/31 \\
2026-09-27T19:15:09 & place & pass & 0/0 \\
2026-09-28T06:05:01 & drc\_routed & pass & 0/0 \\
2026-09-28T06:06:33 & verify & fail & 401/463 \\
2026-09-28T06:21:24 & verify & fail & 402/464 \\
2026-09-28T06:23:51 & verify & fail & 401/463 \\
2026-09-28T06:27:00 & verify & fail & 402/465 \\
2026-09-28T06:28:03 & verify & fail & 382/465 \\
2026-09-28T06:28:18 & place & fail & 15/15 \\
2026-09-28T06:28:30 & place & fail & 15/15 \\
2026-09-28T06:28:43 & place & fail & 15/15 \\
2026-09-28T06:53:50 & verify & fail & 23/458 \\
2026-09-28T06:54:13 & place & fail & 14/14 \\
2026-09-28T07:03:26 & place & fail & 14/14 \\
2026-09-28T07:04:24 & verify & fail & 8/453 \\
2026-09-28T07:25:09 & drc\_routed & pass & 0/0 \\
2026-09-28T07:25:14 & place & pass & 0/14 \\
2026-09-28T07:26:14 & verify & pass & 0/450 \\
2026-09-28T07:31:00 & drc\_routed & pass & 0/0 \\
2026-09-28T07:31:05 & place & pass & 0/14 \\
2026-09-28T07:33:09 & verify & fail & 1/450 \\
2026-09-28T07:33:36 & dfm & pass & 0/4 \\
2026-09-28T07:35:05 & verify & pass & 0/450 \\
2026-09-28T08:11:59 & erc & fail & 60/60 \\
2026-09-28T08:12:07 & drc\_routed & pass & 0/0 \\
2026-09-28T08:12:26 & dfm & pass & 0/4 \\
2026-09-28T08:13:54 & erc & pass & 0/0 \\
2026-09-28T08:14:02 & drc\_routed & pass & 0/0 \\
2026-09-28T08:14:07 & place & pass & 0/14 \\
2026-09-28T08:15:06 & verify & pass & 0/450 \\
2026-09-28T08:15:25 & dfm & pass & 0/4 \\
\bottomrule
\end{longtable}
}
\subsection*{Issues}
\emph{(no entries)}
\subsection*{Human checkpoints}
{\small
\begin{longtable}{lllp{2.6cm}p{6.4cm}}
\toprule
\textbf{Id} & \textbf{Phase} & \textbf{Status} & \textbf{When} & \textbf{Note} \\
\midrule
\endhead
1 & P2 & skipped & 2026-09-27T14:56:54 & Unattended run (owner: decide for yourself): architecture taken as PROVISIONAL delegate answer, not owner approval. 10-28 V, 10 A RMS/20 A pk, 60 V FETs, DRV8300D bootstrap, 3x Kelvin shunts + INA240, STM32G4, LM74500 RP NFET, buck 5V/LDO 3V3/12V boost, 4L. All P2 slots under coverage gap. \\
2 & P4 & not reached &  &  \\
3 & P6 & not reached &  &  \\
4 & P8 & not reached &  &  \\
5 & P10 & not reached &  &  \\
\bottomrule
\end{longtable}
}
\subsection*{Generation cost}
\emph{Not recorded (no reports/cost.json).}
\section{Artifact Index}
{\small
\begin{longtable}{p{9cm}rp{4cm}}
\toprule
\textbf{File} & \textbf{Size} & \textbf{Note} \\
\midrule
\endhead
\texttt{state.json} & 42,624 B &  \\
\texttt{kicad/PCB-0018-A\_bldc-motor-driver.kicad\_pcb} & 3,466,468 B &  \\
\texttt{kicad/PCB-0018-A\_bldc-motor-driver.kicad\_sch} & 48,296 B &  \\
\texttt{fab/PCB-0018-A\_bldc-motor-driver\_gerbers.zip} & 678,768 B &  \\
\texttt{fab/BOM.csv} & 2,500 B &  \\
\texttt{fab/CPL.csv} & 5,233 B &  \\
\texttt{brief/brief.md} & 3,093 B &  \\
\texttt{requirements.md} & 12,843 B &  \\
\texttt{architecture/blocks.md} & 6,707 B &  \\
\texttt{architecture/constraints.json} & 9,248 B &  \\
\texttt{architecture/mcu\_pinmap.md} & 2,269 B &  \\
\texttt{architecture/power\_tree.md} & 3,111 B &  \\
\texttt{architecture/sheets.md} & 2,996 B &  \\
\texttt{architecture/stackup.md} & 1,381 B &  \\
\texttt{reports/PCB-0018-A\_bldc-motor-driver\_bottom.png} & 144,019 B &  \\
\texttt{reports/PCB-0018-A\_bldc-motor-driver\_top.png} & 602,087 B &  \\
\texttt{reports/bom\_cpl.json} & 75,023 B &  \\
\texttt{reports/check\_diffpair.json} & 12,725 B &  \\
\texttt{reports/fab\_export.json} & 2,779 B &  \\
\texttt{reports/place-waivers.json} & 9,956 B &  \\
\texttt{reports/review-schematic.json} & 4,930 B &  \\
\texttt{reports/review-schematic.md} & 4,965 B &  \\
\texttt{reports/schematic.pdf} & 620,917 B &  \\
\texttt{reports/verify-waivers.json} & 209,840 B &  \\
\bottomrule
\end{longtable}
}
\end{document}
